Dans une expérience de Young, une lumière monochromatique de longueur d'onde \qty{\v{lam}}{nm} éclaire deux fentes très fines séparées de \qty{\v{sep}}{mm}. La figure d'interférence est observée sur un écran placé à \qty{\v{L}}{m} des fentes.

\begin{center}
\begin{tikzpicture}[thick]
    \draw (0,-1.8) -- (0,-0.65); \draw (0,-0.55) -- (0,0.55); \draw (0,0.65) -- (0,1.8);
    \draw[very thick] (6,-1.8) -- (6,2.4);
    \node[above] at (0,1.8) {\small fentes};
    \node[above] at (6,2.4) {\small écran};
    \draw[thin,dashed] (0,0) -- (6,0);
    \draw[thin] (0,0.6) -- (6,2);
    \draw[thin] (0,-0.6) -- (6,2);
    \draw[thin,dashed] (0,0) -- (6,2);
    \draw[thin] (1.6,0) arc (0:18.43:1.6);
    \node at (2,0.3) {\small$\theta$};
    \fill (6,2) circle (2pt) node[above right] {\small frange brillante $m = \v{m}$};
    \draw[thin,|-|] (6.4,0) -- node[right] {\small$y$} (6.4,2);
    \draw[thin,|-|] (-0.4,-0.6) -- node[left] {\small\qty{\v{sep}}{mm}} (-0.4,0.6);
    \draw[thin,|-|] (0,-2.2) -- node[below] {\small\qty{\v{L}}{m}} (6,-2.2);
\end{tikzpicture}
\end{center}

\begin{enumerate}[label=\alph*)]
    \item Déterminez l'angle auquel on observe la frange brillante d'ordre $m = \v{m}$.
    \item Déterminez la distance sur l'écran entre le centre de la figure et cette frange brillante.
    \item Déterminez la distance entre deux franges brillantes voisines près du centre de l'écran.
\end{enumerate}
